% Options for packages loaded elsewhere
\PassOptionsToPackage{unicode}{hyperref}
\PassOptionsToPackage{hyphens}{url}
\documentclass[
  english,
  man]{apa6}
\usepackage{xcolor}
\usepackage{amsmath,amssymb}
\setcounter{secnumdepth}{-\maxdimen} % remove section numbering
\usepackage{iftex}
\ifPDFTeX
  \usepackage[T1]{fontenc}
  \usepackage[utf8]{inputenc}
  \usepackage{textcomp} % provide euro and other symbols
\else % if luatex or xetex
  \usepackage{unicode-math} % this also loads fontspec
  \defaultfontfeatures{Scale=MatchLowercase}
  \defaultfontfeatures[\rmfamily]{Ligatures=TeX,Scale=1}
\fi
\usepackage{lmodern}
\ifPDFTeX\else
  % xetex/luatex font selection
\fi
% Use upquote if available, for straight quotes in verbatim environments
\IfFileExists{upquote.sty}{\usepackage{upquote}}{}
\IfFileExists{microtype.sty}{% use microtype if available
  \usepackage[]{microtype}
  \UseMicrotypeSet[protrusion]{basicmath} % disable protrusion for tt fonts
}{}
\makeatletter
\@ifundefined{KOMAClassName}{% if non-KOMA class
  \IfFileExists{parskip.sty}{%
    \usepackage{parskip}
  }{% else
    \setlength{\parindent}{0pt}
    \setlength{\parskip}{6pt plus 2pt minus 1pt}}
}{% if KOMA class
  \KOMAoptions{parskip=half}}
\makeatother
% Make \paragraph and \subparagraph free-standing
\makeatletter
\ifx\paragraph\undefined\else
  \let\oldparagraph\paragraph
  \renewcommand{\paragraph}{
    \@ifstar
      \xxxParagraphStar
      \xxxParagraphNoStar
  }
  \newcommand{\xxxParagraphStar}[1]{\oldparagraph*{#1}\mbox{}}
  \newcommand{\xxxParagraphNoStar}[1]{\oldparagraph{#1}\mbox{}}
\fi
\ifx\subparagraph\undefined\else
  \let\oldsubparagraph\subparagraph
  \renewcommand{\subparagraph}{
    \@ifstar
      \xxxSubParagraphStar
      \xxxSubParagraphNoStar
  }
  \newcommand{\xxxSubParagraphStar}[1]{\oldsubparagraph*{#1}\mbox{}}
  \newcommand{\xxxSubParagraphNoStar}[1]{\oldsubparagraph{#1}\mbox{}}
\fi
\makeatother
\usepackage{graphicx}
\makeatletter
\newsavebox\pandoc@box
\newcommand*\pandocbounded[1]{% scales image to fit in text height/width
  \sbox\pandoc@box{#1}%
  \Gscale@div\@tempa{\textheight}{\dimexpr\ht\pandoc@box+\dp\pandoc@box\relax}%
  \Gscale@div\@tempb{\linewidth}{\wd\pandoc@box}%
  \ifdim\@tempb\p@<\@tempa\p@\let\@tempa\@tempb\fi% select the smaller of both
  \ifdim\@tempa\p@<\p@\scalebox{\@tempa}{\usebox\pandoc@box}%
  \else\usebox{\pandoc@box}%
  \fi%
}
% Set default figure placement to htbp
\def\fps@figure{htbp}
\makeatother
\ifLuaTeX
\usepackage[bidi=basic,shorthands=off]{babel}
\else
\usepackage[bidi=default,shorthands=off]{babel}
\fi
\ifLuaTeX
  \usepackage{selnolig} % disable illegal ligatures
\fi
\setlength{\emergencystretch}{3em} % prevent overfull lines
\providecommand{\tightlist}{%
  \setlength{\itemsep}{0pt}\setlength{\parskip}{0pt}}
% Manuscript styling
\usepackage{upgreek}
\captionsetup{font=singlespacing,justification=justified}

% Table formatting
\usepackage{longtable}
\usepackage{lscape}
% \usepackage[counterclockwise]{rotating}   % Landscape page setup for large tables
\usepackage{multirow}		% Table styling
\usepackage{tabularx}		% Control Column width
\usepackage[flushleft]{threeparttable}	% Allows for three part tables with a specified notes section
\usepackage{threeparttablex}            % Lets threeparttable work with longtable

% Create new environments so endfloat can handle them
% \newenvironment{ltable}
%   {\begin{landscape}\centering\begin{threeparttable}}
%   {\end{threeparttable}\end{landscape}}
\newenvironment{lltable}{\begin{landscape}\centering\begin{ThreePartTable}}{\end{ThreePartTable}\end{landscape}}

% Enables adjusting longtable caption width to table width
% Solution found at http://golatex.de/longtable-mit-caption-so-breit-wie-die-tabelle-t15767.html
\makeatletter
\newcommand\LastLTentrywidth{1em}
\newlength\longtablewidth
\setlength{\longtablewidth}{1in}
\newcommand{\getlongtablewidth}{\begingroup \ifcsname LT@\roman{LT@tables}\endcsname \global\longtablewidth=0pt \renewcommand{\LT@entry}[2]{\global\advance\longtablewidth by ##2\relax\gdef\LastLTentrywidth{##2}}\@nameuse{LT@\roman{LT@tables}} \fi \endgroup}

% \setlength{\parindent}{0.5in}
% \setlength{\parskip}{0pt plus 0pt minus 0pt}

% Overwrite redefinition of paragraph and subparagraph by the default LaTeX template
% See https://github.com/crsh/papaja/issues/292
\makeatletter
\renewcommand{\paragraph}{\@startsection{paragraph}{4}{\parindent}%
  {0\baselineskip \@plus 0.2ex \@minus 0.2ex}%
  {-1em}%
  {\normalfont\normalsize\bfseries\itshape\typesectitle}}

\renewcommand{\subparagraph}[1]{\@startsection{subparagraph}{5}{1em}%
  {0\baselineskip \@plus 0.2ex \@minus 0.2ex}%
  {-\z@\relax}%
  {\normalfont\normalsize\itshape\hspace{\parindent}{#1}\textit{\addperi}}{\relax}}
\makeatother

\makeatletter
\usepackage{etoolbox}
\patchcmd{\maketitle}
  {\section{\normalfont\normalsize\abstractname}}
  {\section*{\normalfont\normalsize\abstractname}}
  {}{\typeout{Failed to patch abstract.}}
\patchcmd{\maketitle}
  {\section{\protect\normalfont{\@title}}}
  {\section*{\protect\normalfont{\@title}}}
  {}{\typeout{Failed to patch title.}}
\makeatother

\usepackage{xpatch}
\makeatletter
\xapptocmd\appendix
  {\xapptocmd\section
    {\addcontentsline{toc}{section}{\appendixname\ifoneappendix\else~\theappendix\fi: #1}}
    {}{\InnerPatchFailed}%
  }
{}{\PatchFailed}
\makeatother
\keywords{behavioral immune system, attentional bias, dot-probe, replication, computational reproducibility\newline\indent Word count: X}
\DeclareDelayedFloatFlavor{ThreePartTable}{table}
\DeclareDelayedFloatFlavor{lltable}{table}
\DeclareDelayedFloatFlavor*{longtable}{table}
\makeatletter
\renewcommand{\efloat@iwrite}[1]{\immediate\expandafter\protected@write\csname efloat@post#1\endcsname{}}
\makeatother
\usepackage{lineno}

\linenumbers
\usepackage{csquotes}
\usepackage{bookmark}
\IfFileExists{xurl.sty}{\usepackage{xurl}}{} % add URL line breaks if available
\urlstyle{same}
\hypersetup{
  pdftitle={Preregistered Direct Replication of Miller \& Maner (2011): Sick Body{,} Vigilant Mind},
  pdfauthor={沈钰涵1},
  pdflang={en-EN},
  pdfkeywords={behavioral immune system, attentional bias, dot-probe, replication, computational reproducibility},
  hidelinks,
  pdfcreator={LaTeX via pandoc}}

\title{Preregistered Direct Replication of Miller \& Maner (2011): Sick Body, Vigilant Mind}
\author{沈钰涵\textsuperscript{1}}
\date{}


\shorttitle{SBVM Replication Report}

\authornote{

Computational reproducibility check of Tybur et al.~(2020).
All analyses were re-run from raw data obtained from the OSF repository.

Correspondence concerning this article should be addressed to 沈钰涵, Nanjing Normal University, School of Psychology. E-mail: \href{mailto:3574425102@qq.com}{\nolinkurl{3574425102@qq.com}}

}

\affiliation{\vspace{0.5cm}\textsuperscript{1} Nanjing Normal University, School of Psychology}

\abstract{%
This report presents a comprehensive computational reproducibility check of Tybur et al.'s (2020) preregistered direct replication of Miller and Maner (2011).
Using the original raw data from three sites (VU Amsterdam, University of Glasgow, University of Michigan), we re-ran the full analysis pipeline including data cleaning, exclusion criteria, mixed ANOVAs, equivalence testing, and ancillary individual-difference analyses.
Results successfully reproduced the key findings: a significant main effect of face type on reaction times (disfigured \textgreater{} typical), a significant main effect of illness recency (recently ill participants showed larger attentional bias), and a non-significant interaction.
Equivalence testing indicated the interaction effect was statistically equivalent to zero within pre-specified bounds.
Ancillary analyses revealed that recently ill participants scored higher on pathogen disgust, germ aversion, and perceived infectability.
}



\begin{document}
\maketitle

\section{Method}\label{method}

\subsection{Participants}\label{participants}

413 participants across three sites completed the dot-probe task:
VU Amsterdam (147), University of Glasgow (121), and University of Michigan (145).
After excluding 9 participants due to excessive errors and 2 participants due to missing questionnaire data,
the final sample consisted of 402 participants
(151 recently ill, 251 not recently ill).

\subsection{Measures}\label{measures}

\subsubsection{Pathogen Disgust Scales}\label{pathogen-disgust-scales}

Internal consistencies (standardized Cronbach's \(\alpha\)) were acceptable to good:
Perceived Infectability \(\alpha = 0.92\), Germ Aversion \(\alpha = 0.76\),
Pathogen Disgust \(\alpha = 0.69\), and Disgust Toward Images \(\alpha = 0.81\).

\section{Results}\label{results}

\subsection{Descriptive Statistics}\label{descriptive-statistics}

\begin{table}

\caption{\label{tab:desc-table}Table 1. Mean reaction times (ms) by face type and illness recency.}
\centering
\begin{tabular}[t]{l|l|r|r|r}
\hline
face\_type & recent\_cold & n & mean & sd\\
\hline
disfigured & recent & 151 & 669.3 & 209.4\\
\hline
disfigured & not recent & 251 & 629.5 & 157.5\\
\hline
typical & recent & 151 & 653.6 & 184.5\\
\hline
typical & not recent & 251 & 622.0 & 148.5\\
\hline
\end{tabular}
\end{table}

Mean reaction times were 644.50 ms
(\(SD = 179.60) for disfigured faces and
633.90 ms
(\)SD = 163.50) for typical faces.

\subsection{Primary Analysis: 2 (Face Type) × 2 (Illness Recency) Mixed ANOVA}\label{primary-analysis-2-face-type-2-illness-recency-mixed-anova}

A 2 (face type: disfigured vs.~typical; within-subjects) × 2 (illness recency: recent vs.~not recent; between-subjects) mixed ANOVA was conducted on mean reaction times from incongruent-location trials.

\subsubsection{Main Effect of Face Type}\label{main-effect-of-face-type}

F(1, 400) = 14.96, \(\eta_{p}^{2}\) = 0.036, 90\% CI = {[}0.009, 0.078{]}, p \textless{} .001

Participants responded more slowly to disfigured faces (\$M = 644.5 \$ ms, \$SD = 179.6 \() than to typical faces (\)M = 633.9 \$ ms, \$SD = 163.5 \$).

\subsubsection{Main Effect of Illness Recency}\label{main-effect-of-illness-recency}

F(1, 400) = 4.23, \(\eta_{p}^{2}\) = 0.010, 90\% CI = {[}0.000, 0.039{]}, p = .040

Recently ill participants showed longer RTs overall compared to not-recently-ill participants.

\subsubsection{Interaction: Face Type × Illness Recency}\label{interaction-face-type-illness-recency}

F(1, 400) = 1.87, \(\eta_{p}^{2}\) = 0.005, 90\% CI = {[}0.000, 0.027{]}, p = .173

The interaction between face type and illness recency was not significant, F(1, 400) = 1.87, \(\eta_{p}^{2}\) = 0.005, 90\% CI = {[}0.000, 0.027{]}, p = .173.

\subsubsection{Simple Effects}\label{simple-effects}

\begin{verbatim}
## **Recently ill group:**  M = 15.71 ms, 95% CI = [4.63, 26.79], p = .006
\end{verbatim}

\begin{verbatim}
## **Not recently ill group:**  M = 7.51 ms, 95% CI = [1.19, 13.83], p = .020
\end{verbatim}

Among recently ill participants, the attentional bias toward disfigured faces (relative to typical faces) was significant, M = 15.71 ms, 95\% CI = {[}4.63, 26.79{]}, p = .006.
Among not-recently-ill participants, this bias was also significant, M = 7.51 ms, 95\% CI = {[}1.19, 13.83{]}, p = .020,
though numerically smaller in magnitude.

\subsection{Equivalence Testing}\label{equivalence-testing}

\begin{verbatim}
## To formally test whether the interaction effect was practically equivalent to zero, we conducted a two-one-sided test (TOST) procedure.
\end{verbatim}

\begin{verbatim}
## The TOST procedure tested whether the difference in attentional bias scores between groups fell within equivalence bounds of $\pm$0.35 Cohen's $d_z$.
\end{verbatim}

\begin{verbatim}
## NHST: $t( 248.4 ) =  -1.27 $,  p = .206
\end{verbatim}

\begin{verbatim}
## TOST lower bound: $t( 248.4 ) =  2.01 $,  p = .023
\end{verbatim}

\begin{verbatim}
## TOST upper bound: $t( 248.4 ) =  -4.55 $,  p < .001
\end{verbatim}

\begin{verbatim}
## The TOST procedure was significant ($p =  0.023 $), indicating that the interaction effect is statistically equivalent to zero within the pre-registered bounds.
\end{verbatim}

\begin{verbatim}
## Cohen's $d_z =  -0.14 $.
\end{verbatim}

\begin{verbatim}
## 90% CI for $d_z$: $[ -0.31 ,  0.04 ]$
\end{verbatim}

\subsection{Continuous Illness Variable Analysis}\label{continuous-illness-variable-analysis}

When using the continuous illness recency variable instead of the dichotomous split, the pattern remained consistent:

Face type: F(1, 400) = 13.21, \(\eta_{p}^{2}\) = 0.032, 90\% CI = {[}0.007, 0.073{]}, p \textless{} .001

Recency (continuous): F(1, 400) = 1.06, \(\eta_{p}^{2}\) = 0.003, 90\% CI = {[}0.000, 0.022{]}, p = .303

Interaction: F(1, 400) = 0.16, \(\eta_{p}^{2}\) \textless{} .001, 90\% CI = {[}0.000, 0.013{]}, p = .691

\subsection{Ancillary Analyses: Individual Differences}\label{ancillary-analyses-individual-differences}

\subsubsection{Illness Recency and Disgust Sensitivity}\label{illness-recency-and-disgust-sensitivity}

Recently ill participants scored significantly higher on:

\begin{itemize}
\tightlist
\item
  Perceived Infectability: t(311.3) = 2.78, p = .006 , r = .16, CI = {[}.05, .26{]}, p = .006
\item
  Germ Aversion: t(338.6) = 2.09, p = .037 , r = .11, CI = {[}.01, .22{]}, p = .037
\item
  Pathogen Disgust: t(299.8) = 2.08, p = .038 , r = .12, CI = {[}.01, .23{]}, p = .038
\item
  Disgust Toward Images: t(308.0) = 0.92, p = .358 , r = .05, CI = {[}-.06, .16{]}, p = .358
\end{itemize}

\subsubsection{Correlations Between Recency and Individual Differences}\label{correlations-between-recency-and-individual-differences}

Continuous illness recency correlated significantly with:

\begin{itemize}
\tightlist
\item
  Perceived Infectability: r = .26, CI = {[}.17, .35{]}, p \textless{} .001
\item
  Germ Aversion: r = .08, CI = {[}-.02, .17{]}, p = .127
\item
  Pathogen Disgust: r = .09, CI = {[}-.00, .19{]}, p = .062
\item
  Disgust Toward Images: r = .10, CI = {[}.00, .19{]}, p = .047
\end{itemize}

\subsubsection{Face-Type Correlation}\label{face-type-correlation}

\begin{verbatim}
## Reaction times for disfigured and typical faces were highly correlated,  r = .95 ,  95% CI = [0.94, 0.96] ,  p < .001 .
\end{verbatim}

\subsection{Robustness Checks}\label{robustness-checks}

\subsubsection{Winsorized (Trimmed) Data}\label{winsorized-trimmed-data}

When using winsorized RTs (replacing outliers \$\pm\$3 SD with boundary values) instead of exclusion:

\textbf{Categorical ANOVA:}
Face type: \(F(1, 400) = 22.21\), \(\eta_{p}^{2} = .053\), \(p < .001\)
Recent cold: \(F(1, 400) = 4.67\), \(\eta_{p}^{2} = .012\), \(p = .031\)
Interaction: \(F(1, 400) = 2.63\), \(\eta_{p}^{2} = .007\), \(p = .105\)

\subsubsection{Per-Site Analyses}\label{per-site-analyses}

\textbf{Michigan (MI):}
Face type: \(F(1, 138) = 7.26\), \(\eta_{p}^{2} = .050\), \(p = .008\)
Recent cold: \(F(1, 138) = 3.36\), \(\eta_{p}^{2} = .024\), \(p = .069\)
Interaction: \(F(1, 138) = 0.53\), \(\eta_{p}^{2} = .004\), \(p = .467\)

\textbf{VU Amsterdam (VU):}
Face type: \(F(1, 142) = 1.26\), \(\eta_{p}^{2} = .009\), \(p = .263\)
Recent cold: \(F(1, 142) = 0.06\), \(\eta_{p}^{2} < .001\), \(p = .809\)
Interaction: \(F(1, 142) = 0.01\), \(\eta_{p}^{2} < .001\), \(p = .908\)

\textbf{Glasgow (GL):}
Face type: \(F(1, 116) = 6.39\), \(\eta_{p}^{2} = .052\), \(p = .013\)
Recent cold: \(F(1, 116) = 1.24\), \(\eta_{p}^{2} = .011\), \(p = .269\)
Interaction: \(F(1, 116) = 1.26\), \(\eta_{p}^{2} = .011\), \(p = .263\)

\subsubsection{ANCOVA (Controlling for Typical RT)}\label{ancova-controlling-for-typical-rt}

\textbf{Categorical ANCOVA} (controlling for typical-face RT):
Typical RT: \(F(1, 399) = 3396.32\), \(\eta_{p}^{2} = .895\), \(p < .001\)
Recent cold: \(F(1, 399) = 1.36\), \(\eta_{p}^{2} = .003\), \(p = .244\)

\textbf{Continuous ANCOVA:}
Typical RT: \(F(1, 399) = 3418.14\), \(\eta_{p}^{2} = .895\), \(p < .001\)
Recency (z): \(F(1, 399) = 0.08\), \(\eta_{p}^{2} < .001\), \(p = .773\)

\subsubsection{Linear Mixed-Effects Models}\label{linear-mixed-effects-models}

\textbf{Categorical LMEM} (random intercepts for site, participant, and stimulus):
Recent cold (\(\beta = 32.53\), \(SE = 16.52\), \(p = .050\))
Face type (\(\beta = -12.59\), \(SE = 11.20\), \(p = .268\))
Interaction (\(\beta = -8.26\), \(SE = 6.12\), \(p = .178\))

\textbf{Continuous LMEM:}
Recency (\(\beta = 11.29\), \(SE = 8.01\), \(p = .160\))
Face type (\(\beta = -11.56\), \(SE = 11.17\), \(p = .308\))
Interaction (\(\beta = -1.36\), \(SE = 2.97\), \(p = .647\))

\subsubsection{HEXACO Personality Factors}\label{hexaco-personality-factors}

No significant differences in HEXACO traits were found between recently and not-recently ill participants (all \(p\)s \(> .05\)).

\section{Discussion}\label{discussion}

The present computational reproducibility check successfully reproduced the key findings of Tybur et al.~(2020).
Participants consistently responded more slowly to disfigured faces compared to typical faces on incongruent-location trials,
indicating an attentional bias toward cues of disease.
Although recently ill participants showed numerically larger bias scores, the interaction between face type and illness recency
was not statistically significant, and equivalence testing confirmed the interaction was practically equivalent to zero
within pre-registered bounds of \$\textbackslash pm\$0.35 \(d_z\).

Ancillary analyses replicated the finding that recently ill participants reported higher levels of pathogen disgust sensitivity,
germ aversion, and perceived infectability, consistent with the behavioral immune system framework.
Robustness checks using winsorized data, per-site analyses, ANCOVAs, and linear mixed-effects models
all converged on the same pattern of results, supporting the robustness of the findings.

\section{References}\label{references}

\protect\phantomsection\label{refs}


\end{document}
